\documentclass[11pt,a4paper]{article}
\usepackage[top=1.7cm,bottom=1.5cm,left=1.5cm,right=1.5cm,headsep=8pt,footskip=16pt]{geometry}
\usepackage{enumitem}
\usepackage{amssymb}
\usepackage{array}
\usepackage[table]{xcolor}
\usepackage{fancyhdr}
\usepackage{lastpage}
\usepackage{multicol}
\usepackage{polyglossia}
\setdefaultlanguage[calendar=gregorian, numerals=maghrib]{arabic}
\setotherlanguage{english}
\IfFontExistsTF{Amiri}{%
  \newfontfamily\arabicfont[Script=Arabic, Scale=1.08]{Amiri}
}{%
  \newfontfamily\arabicfont[Script=Arabic, Scale=1.08, Path=../fonts/, Extension=.ttf,
    UprightFont=*-Regular, BoldFont=*-Bold, ItalicFont=*-Italic,
    BoldItalicFont=*-BoldItalic]{Amiri}
}
\IfFontExistsTF{AlBattar}{%
  \newfontfamily\khatt[Script=Arabic, Scale=1.15]{AlBattar}
}{%
  \newfontfamily\khatt[Script=Arabic, Scale=1.15, Path=../fonts/,
    BoldFont=ae_AlBattar.ttf]{ae_AlBattar.ttf}
}
\renewcommand{\headrulewidth}{0.8pt}
\renewcommand{\footrulewidth}{0.8pt}
\fancyhf{}
\rhead{\khatt المدرسة العليا للأساتذة ببوزريعة}
\lhead{\khatt ملحقة بومرداس}
\lfoot{\khatt الأستاذ: د. بن عبد الرزاق يوسف}
\cfoot{\thepage/\pageref{LastPage}}
\pagestyle{fancy}
\parindent 0pt
\setlength{\columnsep}{0.6cm}
\setlength{\columnseprule}{0.4pt}
\setlist[enumerate,1]{label=\textbf{\arabic*.}, leftmargin=*, itemsep=5pt, topsep=2pt}
\begin{document}

\linespread{1.08}\selectfont
\begin{center}
{\LARGE\khatt\textbf{الاختبار الثامن في الذكاء الاصطناعي}}\\[4pt]
{\normalsize وضعيات إشكالية ــ أسئلة متعددة الاختيارات ــ الفصول من 1 إلى 6}
\end{center}
\vspace{-2pt}\rule{\textwidth}{0.6pt}\\[2pt]
\textbf{التعليمة:} يتضمن الاختبار خمس وضعيات إشكالية. اقرأ كل وضعية بعناية، ثم اختر التحليل أو التصرف الأنسب
من بين الاختيارات الخمسة؛ لكل وضعية إجابة صحيحة واحدة فقط.
اكتب \textbf{حرف} الإجابة الصحيحة في جدول الإجابات في آخر الاختبار.
لكل وضعية 4 نقاط (المجموع 20 نقطة).
\vspace{4pt}

\noindent
{\large\khatt\textbf{الوضعية 1}}\hfill{\small(الفصل 1)}\par\nopagebreak\vspace{3pt}
\setlength{\fboxsep}{7pt}\colorbox{gray!10}{\parbox{\dimexpr\linewidth-14pt}{سأل أستاذ أداة ذكاء اصطناعي عن محتوى منهاج اللغة العربية الجديد للسنة الأولى متوسط، الصادر حديثًا، فقدّمت إجابة مفصّلة واثقة، لكنها تطابق المنهاج القديم الملغى. ويرى زميله أن الأداة «لا تكذب»، وأن الخطأ في فهم الأستاذ للمنهاج. ويستعد لإعداد التدرجات السنوية بناءً على هذه المعلومات. ما التفسير الأدق؟}}\par\nopagebreak
\begin{itemize}[leftmargin=1.0cm, labelsep=6pt, itemsep=3pt, topsep=4pt, parsep=0pt]
  \item[\textbf{أ)}] الأداة تعمّدت تقديم المنهاج القديم لأنها لا تثق في التغييرات الحديثة، والنماذج مصممة لتفضيل المعلومات المستقرة، فتصبح إجابتها مقبولة تربويًا ولو خالفت المنهاج الجديد.
  \item[\textbf{ب)}] معرفة النموذج محدودة بتاريخ بيانات تدريبه، فإذا صدر المنهاج بعده قدّم ما تعلّمه سابقًا بثقة، لأن الثقة في الصياغة لا تعكس حداثة المعلومة؛ فالمرجع الوثائق الرسمية للوزارة.
  \item[\textbf{ج)}] الأداة تعرف المنهاج الجديد لكنها تفضّل القديم لأنه أكثر انتشارًا في الإنترنت، فيكفي أن يطلب منها «المنهاج الجديد تحديدًا» لتعطيه المحتوى الصحيح دون الرجوع إلى أي وثيقة رسمية.
  \item[\textbf{د)}] معرفة النموذج محدودة بتاريخ بياناته، والحل أن يسأل الأداة نفسها عن تاريخ صدور المنهاج الجديد، فإن ذكرت تاريخًا سابقًا لبياناتها اعتمد إجابتها، وإلا بحث في وثائق الوزارة.
  \item[\textbf{هـ)}] زميله محق، فالأداة لا تقدّم إلا معلومات صحيحة، والخطأ في المنهاج الجديد نفسه الذي لم يُنشر رسميًا بعد، ولذلك يُعتمد ما قدّمته الأداة حتى تصدر الوزارة توضيحًا بشأنه.
\end{itemize}
\par\vspace{6pt}
\noindent
{\large\khatt\textbf{الوضعية 2}}\hfill{\small(الفصل 2)}\par\nopagebreak\vspace{3pt}
\setlength{\fboxsep}{7pt}\colorbox{gray!10}{\parbox{\dimexpr\linewidth-14pt}{جرّب أساتذة أداة تحويل الكلام إلى نص على قراءات تلاميذ من مناطق مختلفة، فكانت دقيقة مع النطق المشرقي، وضعيفة مع النطق المغاربي، خاصة في بعض الأصوات والوقفات. واقترح كل أستاذ تفسيرًا لذلك. ويرغب الأساتذة في تقرير يُرفع إلى الإدارة حول صلاحية الأداة. أي تفسير هو الأقرب إلى طريقة تعلّم النماذج؟}}\par\nopagebreak
\begin{itemize}[leftmargin=1.0cm, labelsep=6pt, itemsep=3pt, topsep=4pt, parsep=0pt]
  \item[\textbf{أ)}] بيانات تدريب الأداة يغلب عليها على الأرجح النطق المشرقي، فتعلّمت أنماطه أكثر من غيره؛ وتحسين أدائها يحتاج إلى بيانات صوتية كافية تمثّل النطق المغاربي، لأن النموذج يتقن ما رآه كثيرًا.
  \item[\textbf{ب)}] بيانات التدريب يغلب عليها النطق المشرقي، والحل أن يكتب الأستاذ في كل طلب «المتحدث مغاربي»، لأن هذه العبارة تجعل النموذج يعيد تدريب نفسه فورًا على النطق المغاربي، دون حاجة إلى جمع بيانات جديدة.
  \item[\textbf{ج)}] الأداة تعرف موقع المستخدم الجغرافي، فتخفض دقتها عمدًا في المناطق التي لا تدعمها الشركة تجاريًا، ولا علاقة لذلك ببيانات التدريب أو بخصائص الأصوات المنطوقة.
  \item[\textbf{د)}] النطق المغاربي أقل فصاحة من المشرقي، ولذلك تعجز الأداة عنه، والحل تدريب التلاميذ على النطق المشرقي، لأن الأداة تقيس الفصاحة بدقة، وضعف أدائها مؤشر على ضعف نطق التلاميذ، ويُعدّ ذلك تدريبًا مفيدًا لهم على الفصاحة.
  \item[\textbf{هـ)}] الأصوات المغاربية لا يمكن تمثيلها رقميًا لأنها لا تُكتب بالحروف العربية المعيارية، ولذلك يستحيل على أي نموذج تعلّمها مهما كانت كمية البيانات المتوفرة أو جودتها.
\end{itemize}
\par\vspace{6pt}
\noindent
{\large\khatt\textbf{الوضعية 3}}\hfill{\small(الفصل 4)}\par\nopagebreak\vspace{3pt}
\setlength{\fboxsep}{7pt}\colorbox{gray!10}{\parbox{\dimexpr\linewidth-14pt}{تريد أستاذة تحويل صورة جماعية لتلاميذ قسمها إلى رسوم كرتونية بأداة ذكاء اصطناعي مجانية، ثم نشرها في مجلة المؤسسة وعلى صفحتها في مواقع التواصل الاجتماعي. التلاميذ قُصَّر، ولم يُستشر أولياؤهم، ولم تطّلع على شروط استعمال الأداة. وترى أن النشاط سيفرح التلاميذ ويشجعهم على المشاركة في المجلة. ما الإجراء الأسلم؟}}\par\nopagebreak
\begin{itemize}[leftmargin=1.0cm, labelsep=6pt, itemsep=3pt, topsep=4pt, parsep=0pt]
  \item[\textbf{أ)}] تنشر الرسوم في المجلة الورقية فقط دون موافقة، لأن مشكلة الخصوصية تقتصر على مواقع التواصل، أما رفع الصورة الأصلية إلى الأداة فلا إشكال فيه، لأن الأدوات لا تحتفظ بالصور المرفوعة.
  \item[\textbf{ب)}] تطلب موافقة الأولياء على النشر في مواقع التواصل فقط، وترفع الصورة إلى الأداة دون التحقق من شروطها، لأن الموافقة على النشر تشمل تلقائيًا كل معالجة رقمية تسبقه مهما كانت.
  \item[\textbf{ج)}] لا حاجة إلى موافقة، لأن الرسوم الكرتونية لا تُظهر الوجوه الحقيقية، فتتحول الصورة إلى عمل فني جديد لا يتضمن أي معطى شخصي، حتى لو بقي التلاميذ معروفين لزملائهم وأوليائهم.
  \item[\textbf{د)}] تكتفي بموافقة التلاميذ أنفسهم شفهيًا في القسم، لأنهم أصحاب الصورة، وموافقتهم كافية قانونيًا وأخلاقيًا، ولا يلزم إشراك الأولياء أو المؤسسة في قرار يخص نشاطًا تربويًا داخليًا.
  \item[\textbf{هـ)}] تحصل أولًا على موافقة مكتوبة من الأولياء وفق تعليمات المؤسسة، وتتحقق من شروط الأداة في تخزين الصور، أو تكتفي برسوم لشخصيات متخيَّلة، لأن صور القُصّر معطيات شخصية حساسة.
\end{itemize}
\par\vspace{6pt}
\noindent
{\large\khatt\textbf{الوضعية 4}}\hfill{\small(الفصل 5)}\par\nopagebreak\vspace{3pt}
\setlength{\fboxsep}{7pt}\colorbox{gray!10}{\parbox{\dimexpr\linewidth-14pt}{طلب أستاذ من أداة تلخيص مقالين متعارضين حول استعمال العامية في تدريس العربية، لعرضهما في نقاش مع الطلبة. جاء ملخص المقال الأول مفصّلًا بحججه وأمثلته، وملخص الثاني في سطرين مع وصفه بأنه «رأي تقليدي متجاوز». والنقاش بعد ساعة. وأحد الطلبة يميل مسبقًا إلى الرأي الأول. ما التقدير الأدق؟}}\par\nopagebreak
\begin{itemize}[leftmargin=1.0cm, labelsep=6pt, itemsep=3pt, topsep=4pt, parsep=0pt]
  \item[\textbf{أ)}] يعرض الملخصين ويترك للطلبة اكتشاف الانحياز بأنفسهم دون تنبيه، لأن اكتشاف الانحياز مهارة نقدية، والعرض غير المتوازن أنفع تربويًا من العرض المتوازن في كل نقاش.
  \item[\textbf{ب)}] الملخص سليم، لأن الأداة محايدة بطبيعتها، وقِصَر ملخص المقال الثاني دليل على ضعف حججه، ووصفه بالتقليدي توصيف موضوعي توصلت إليه الأداة بعد تحليل الحجج في المقالين.
  \item[\textbf{ج)}] الخلل الوحيد هو وصف «رأي تقليدي»، فيحذفه ويعرض الملخصين كما هما، لأن اختلاف الطول بين الملخصين أمر طبيعي لا يؤثر في عرض الرأيين بإنصاف أمام الطلبة في النقاش، فيكون العرض أسرع وأوضح للطلبة.
  \item[\textbf{د)}] الملخص غير متوازن، إذ أضعف أحد الرأيين وأضاف إليه حكمًا قيميًا ليس في المقال؛ فيطلب تلخيص كل مقال بالطول والبنية نفسيهما دون أحكام، ثم يقابل الملخصين بالمقالين قبل عرضهما.
  \item[\textbf{هـ)}] يطلب ملخصًا جديدًا يعطي المقال الثاني طولًا أكبر من الأول لتعويض الحيف، ثم يعرضه دون مقابلته بالمقالين، لأن التعويض في الطول يعيد التوازن بين الرأيين بصفة تلقائية، وهذا يكفي لإنصاف الرأي الثاني.
\end{itemize}
\par\vspace{6pt}
\noindent
{\large\khatt\textbf{الوضعية 5}}\hfill{\small(الفصل 6)}\par\nopagebreak\vspace{3pt}
\setlength{\fboxsep}{7pt}\colorbox{gray!10}{\parbox{\dimexpr\linewidth-14pt}{في قسم أستاذة تلميذ كفيف، ودرس الغد يعتمد على لوحة فنية يصفها التلاميذ ويحللون عناصرها. تفكر في استعمال أداة متعددة الأنماط لوصف اللوحة وصفًا نصيًا، ثم تحويله إلى صوت. وتريد أن يشارك التلميذ في النشاط كزملائه. ولم يسبق لها أن استعملت هذه الأدوات مع هذا التلميذ. ما الاستعمال الأنسب لتحقيق الإدماج؟}}\par\nopagebreak
\begin{itemize}[leftmargin=1.0cm, labelsep=6pt, itemsep=3pt, topsep=4pt, parsep=0pt]
  \item[\textbf{أ)}] تعطي التلميذ الوصف المكتوب العادي ليقرأه في القسم، لأن الوصف النصي كافٍ وحده، ولا حاجة إلى تحويله إلى صوت أو إلى أي وسيلة أخرى تناسب حالته البصرية، كما يحدث في أغلب الأقسام.
  \item[\textbf{ب)}] تولّد الوصف وتحوّله إلى صوت دون مراجعة، لأن الأدوات متعددة الأنماط تصف اللوحات الفنية بدقة تامة، ومراجعة الأستاذة قد تُدخل فيه انطباعات شخصية لا حاجة إليها في الدرس، فتوفّر وقتًا ثمينًا قبل الحصة.
  \item[\textbf{ج)}] تولّد وصفًا مفصّلًا للوحة وتراجعه لتصحيح ما أخطأت فيه الأداة وإضافة ما يخدم هدف الدرس، ثم تحوّله إلى صوت أو تقرؤه، فيشارك التلميذ في التحليل انطلاقًا من الوصف المسموع.
  \item[\textbf{د)}] تطلب من الأداة توليد لوحة جديدة أبسط من اللوحة الأصلية، لأن اللوحة البسيطة أسهل وصفًا، ثم تعرضها على التلميذ الكفيف ليصفها بنفسه دون أي مساعدة صوتية.
  \item[\textbf{هـ)}] تعفي التلميذ الكفيف من النشاط كليًا، لأن أدوات الوصف لا تستطيع نقل الأعمال الفنية، والأنسب أن يُكلَّف بنشاط مختلف تمامًا لا علاقة له بموضوع الدرس ولا بزملائه.
\end{itemize}
\par\vspace{6pt}
\begin{center}
{\khatt\textbf{جدول الإجابات}}\\[3pt]
\renewcommand{\arraystretch}{1.6}
\begin{tabular}{|>{\centering\arraybackslash}m{1.9cm}|*{5}{>{\centering\arraybackslash}m{1.2cm}|}}
\hline
\textbf{رقم السؤال} & \textbf{1} & \textbf{2} & \textbf{3} & \textbf{4} & \textbf{5} \\ \hline
\textbf{الإجابة} &  &  &  &  &  \\[8pt] \hline
\end{tabular}

\vspace{8pt}
\renewcommand{\arraystretch}{2.2}
\begin{tabular}{|>{\centering\arraybackslash}m{7.5cm}|>{\centering\arraybackslash}m{3.2cm}|>{\centering\arraybackslash}m{3.2cm}|}
\hline
\textbf{الاسم واللقب} & \textbf{الفوج} & \textbf{العلامة} \\ \hline
 & & \hfill /20 \\ \hline
\end{tabular}
\end{center}
\end{document}
